\section{Descripción y alcance de la solución}

La Parte~II propone una aplicación, programa o bot que permita registrar un número telefónico, solicitar la ubicación del dispositivo asociado y visualizarla en un mapa, en tiempo real o por intervalos. La solución debe operar con consentimiento informado del usuario y con datos anonimizados; el objetivo es académico y no comercial.

\begin{table}[H]
	\caption{Niveles de viabilidad del rastreo según el servicio}
	\label{tab:niveles-rastreo}
	\centering
	\footnotesize
	\renewcommand{\arraystretch}{1.25}
	\begin{tabularx}{\textwidth}{@{}>{\raggedright\arraybackslash}p{0.18\textwidth} >{\raggedright\arraybackslash}p{0.20\textwidth} Y >{\raggedright\arraybackslash}p{0.20\textwidth}@{}}
		\toprule
		\textbf{Servicio} & \textbf{Fuente de ubicación} & \textbf{Viabilidad y límites} & \textbf{Requisito de consentimiento} \\
		\midrule
		WhatsApp & Ubicación compartida manualmente & Limitada e indirecta: la aplicación no expone la ubicación en tiempo real ni permite rastrear a un contacto sin su acción. & El usuario debe compartir su ubicación de forma explícita. \\
		Telegram & Bot y ubicación enviada por el usuario & Un bot puede solicitar la ubicación mediante botones, pero no puede rastrear automáticamente sin que el usuario la envíe. & El usuario debe iniciar el chat y autorizar el envío. \\
		Celular directamente & GPS del dispositivo & Requiere una aplicación instalada en el dispositivo objetivo y permisos de localización. & Consentimiento expreso durante la instalación y uso. \\
		\bottomrule
	\end{tabularx}
\end{table}

\subsection{Alcance propuesto}

\begin{itemize}[leftmargin=*, topsep=2pt, itemsep=2pt]
	\item Prototipo funcional con registro por número telefónico, autenticación y visualización de ubicaciones en un mapa.
	\item Actualización periódica de la ubicación, con intervalo configurable (por ejemplo, cada 10 minutos).
	\item Historial de ubicaciones y, como funcionalidad opcional, alertas por geofencing.
	\item Documentación técnica, modelo de datos, API, plan de pruebas y análisis de riesgos.
	\item Demostración con datos sintéticos o de laboratorio; no se utilizarán personas reales ni dispositivos de terceros.
\end{itemize}

\subsection{Fuera de alcance}

\begin{itemize}[leftmargin=*, topsep=2pt, itemsep=2pt]
	\item Rastrear personas sin consentimiento, interceptar comunicaciones o vulnerar la seguridad de WhatsApp, Telegram o de un operador móvil.
	\item Utilizar APIs no oficiales, ingeniería inversa o técnicas que incumplan los términos de servicio de las plataformas.
	\item Almacenar datos personales reales, credenciales o ubicaciones de terceros en el repositorio o en el informe.
	\item Automatizar bloqueos, decisiones disciplinarias o acciones sobre personas a partir de una puntuación de riesgo.
\end{itemize}

\begin{nota}{Principio de diseño}
La solución debe pedir la ubicación, no tomarla por la fuerza. El usuario objetivo debe conocer quién recibe sus datos, para qué se usan, durante cuánto tiempo se conservan y cómo puede revocar el consentimiento.
\end{nota}
